\documentclass[spanish,letterpaper,oneside]{article}
\def\documenttitle {S6. Foro de Derechos Humanos}
\def\documentsubtitle {Dignidad ontológica, derechos y Estado de Derecho}
\def\documentsubject {Actividad 6 - Teoría del Estado y Constitución}
\def\documentauthor {Martín Jonathan De la Cruz Muñoz}
\def\coursename {Teoría del Estado y Constitución}
\def\coursecode {020}
\def\hfwidthwrap {true}
\def\hfwidthtitle {0.52}
\def\hfwidthcourse {0.42}
\def\universityname {Universidad Abierta y a Distancia de México}
\def\universityfaculty {Licenciatura en Derecho}
\def\universitydepartment {Teoría del Estado y Constitución}
\def\universitydepartmentimage {departamentos/UnADM}
\def\universitydepartmentimagecfg {height=1.57cm}
\def\universitylocation {Roma Norte, Ciudad de México}
\def\authortable {
  \begin{tabular}{ll}
    \textbf{Alumno:} & Martín Jonathan De la Cruz Muñoz \\
    Matrícula: & ES261120204 \\
    Figura docente: & Verónica Ortiz López \\
    Grupo: & 020 \\
    Semestre/Bloque: & 2 / 2 \\
    Semana: & 6 \\
    Estado: & Preparado (pendiente validación de envío) \\
    \multicolumn{2}{l}{Fecha: 8 de noviembre de 2026} \\
    \multicolumn{2}{l}{\universitylocation}
  \end{tabular}
}

% Wrapper local para apuntar a la plantilla maestra de AulaTeX
\input{C:/Users/Sysx/Documents/AulaTeX-Academico/aulaTEX/base/Plantilla-Informe/template.tex}

\def\indexstyle{}
\def\nameltcont{Tabla de contenido}
\setcounter{tocdepth}{2}
\usepackage{helvet}
\renewcommand{\familydefault}{\sfdefault}
\usepackage[most]{tcolorbox}
\usepackage{attachfile}
\usepackage{enumitem}
\usepackage{tikz}
\usetikzlibrary{babel}
\definecolor{unadmGreenDark}{HTML}{174A3A}
\definecolor{unadmGreen}{HTML}{5F8F3A}
\definecolor{unadmGold}{HTML}{B88A2A}
\definecolor{unadmPaper}{HTML}{F6F7F2}
\definecolor{unadmInk}{HTML}{1F2A24}

\attachfilesetup{color={0.373 0.561 0.227},print=false}

\newcommand{\foroCopyButton}[1]{%
  \textattachfile[appearance={},description={Participación íntegra para el foro},mimetype=text/plain]{#1}{%
    \colorbox{unadmGreen}{\small\textcolor{white}{\,Copiar participación\,}}%
  }%
}

\newtcolorbox{forobox}[1][]{%
  enhanced, breakable,
  colback=unadmPaper, colframe=unadmGreen,
  boxrule=0pt, leftrule=3pt, arc=1.2mm,
  left=4mm, right=4mm, top=3mm, bottom=3mm,
  fonttitle=\bfseries\color{white}, coltitle=white, colbacktitle=unadmGreenDark,
  attach boxed title to top left={xshift=4mm,yshift=-2mm},
  boxed title style={colframe=unadmGreenDark,arc=0.6mm},
  #1
}

\setcitestyle{authoryear,open={(},close={)},citesep={;}}
\def\pagemargintop{2.3}
\def\pagemarginbottom{2.3}
\def\predocpageromannumber{true}
\def\predocresetpagenumber{true}

\newcommand{\insertcoverwatermark}{%
  \AddToShipoutPictureBG*{%
    \begin{tikzpicture}[remember picture,overlay]
      \node[opacity=0.13,inner sep=0pt,yshift=-0.35cm]
        at (current page.center) {\includegraphics[width=0.72\paperwidth]{img/departamentos/UnADM.pdf}};
    \end{tikzpicture}%
  }%
}

\setlength{\parindent}{0pt}
\setlength{\parskip}{3pt}

\begin{document}
\insertcoverwatermark
\templatePortrait
\templatePagecfg
\hypersetup{urlcolor=unadmGreenDark,citecolor=black,linkcolor=unadmGreenDark,linktoc=all}
\phantomsection\addcontentsline{toc}{section}{Resumen}
\begin{abstractd}
El presente reporte académico documenta de forma sistemática y rigurosa la aportación académica elaborada para la actividad \emph{S6. Foro de Derechos Humanos} de la asignatura Teoría del Estado y Constitución en la Universidad Abierta y a Distancia de México (UnADM). En el marco del Estado constitucional y democrático de Derecho, se abordan las tres interrogantes directrices planteadas en la planeación docente: la conceptualización ontológica y garantista de los derechos humanos, la discusión en torno a la indivisibilidad frente a la jerarquía valorativa axiológica encabezada por la dignidad humana, y la relación sustantiva de co-originariedad y límite al poder político entre derechos fundamentales y Estado de Derecho. El documento incluye el texto íntegro copiable de la intervención, los fundamentos jurisprudenciales de la Suprema Corte de Justicia de la Nación y la doctrina constitucional aplicable, manteniéndose en estado pendiente de validación por el usuario antes de su publicación en el aula virtual.

\footnote{Declaración sobre el uso ético de inteligencia artificial: Se utilizó asistencia computacional como herramienta de verificación ortotipográfica y soporte editorial conforme a los lineamientos institucionales de la UnADM. La autoría intelectual, el análisis conceptual, la argumentación jurídica y las conclusiones son de la exclusiva responsabilidad del estudiante.}

\textbf{Palabras clave:} derechos humanos; dignidad humana; Estado de Derecho; control de regularidad constitucional; indivisibilidad; interdependencia; Ferrajoli; Bobbio.
\end{abstractd}

\renewcommand{\cftsecleader}{\cftdotfill{\cftdotsep}}
\renewcommand{\cftsubsecleader}{\cftdotfill{\cftdotsep}}
\templateIndex
\templateFinalcfg
\renewcommand{\footrulewidth}{0pt}
\def\documenttitlehf {S6. Foro de Derechos Humanos}
\hypersetup{urlcolor=unadmGreenDark,citecolor=black,linkcolor=unadmGreenDark}
\fontsize{12}{13.8}\selectfont
\sectionfont{\color{unadmGreenDark}\fontsize{14}{16}\selectfont\bfseries}
\subsectionfont{\color{unadmGreenDark}\fontsize{11.3}{13}\selectfont\bfseries}
\titlespacing*{\section}{0pt}{7pt}{4pt}
\titlespacing*{\subsection}{0pt}{5pt}{3pt}

\section{Identificación institucional y estado de la actividad}
\hypertarget{cuerpo-inicio}{}
El presente producto editorial corresponde a la actividad obligatoria de la Semana 6 del calendario escolar oficial de la Licenciatura en Derecho de la UnADM (Semestre 2, Bloque 2), denominada \textbf{S6. Foro de Derechos Humanos} (\url{https://aulavirtual.unadmexico.mx/mod/forum/view.php?id=221688}).

\textbf{Estado procesal:} Borrador preparado para revisión editorial y académica. De conformidad con las directivas de soberanía del usuario del ecosistema AulaTeX, la autorización para redactar y maquetar la intervención no equivale a la autorización para su envío o publicación automática; el artefacto queda depositado en la carpeta de entregas en espera de la validación explícita del usuario. Asimismo, en estricto apego a la honestidad académica, no se inventan réplicas simuladas ni interacciones artificiales con compañeros de grupo; tales intercambios se incorporarán una vez que el foro sea validado y abierto en la plataforma.

\section{Desarrollo temático: Dignidad ontológica, derechos y Estado de Derecho}

\begin{forobox}[title={Participación para el Foro de Derechos Humanos -- Texto para publicación}]
\textbf{Título de la aportación:} Dignidad ontológica, interdependencia de derechos fundamentales y legitimidad sustantiva del Estado constitucional de Derecho. \hfill \foroCopyButton{foro-participacion-Actividad-6.txt}

\bigskip
\textbf{Estimada docente titular y compañeros de grupo:}

Comparto mi aportación reflexiva y dogmática a las preguntas detonadoras planteadas para esta Semana 6:

\medskip
\textbf{1. ¿Qué son los derechos humanos?}

Desde una perspectiva iusfundamental y garantista contemporánea \citep{ferrajoli2001,bobbio1991}, los derechos humanos son aquellas facultades, libertades y prerrogativas inherentes a la condición ontológica de toda persona humana por el mero hecho de existir, cuya titularidad universal no depende de concesiones graciosas del poder estatal, sino que antecede y limita a la autoridad política.

En el ordenamiento jurídico mexicano, a partir de la paradigmática reforma constitucional de 10 de junio de 2011 al artículo 1.º de la CPEUM, los derechos humanos han dejado de concebirse como simples ``garantías individuales otorgadas'' para ser formalmente \textbf{reconocidos} e integrados en un bloque de regularidad constitucional supremo (compuesto por la Constitución y los tratados internacionales de los que el Estado mexicano sea parte, conforme a la Jurisprudencia P./J. 20/2014 de la SCJN). Asimismo, están regidos de manera inexcusable por los principios de \textbf{universalidad, interdependencia, indivisibilidad y progresividad}, imponiendo a todas las autoridades del Estado la obligación categórica de promoverlos, respetarlos, protegerlos y garantizarlos \citep{scjnParametro2014}.

\medskip
\textbf{2. ¿Cuál de los derechos humanos es el más importante?}

Desde el punto de vista de la dogmática constitucional estricta y del principio de \textbf{indivisibilidad e interdependencia} (CPEUM, art. 1.º, párrafo tercero), no existe una jerarquía cuantitativa entre derechos humanos: la privación o menoscabo de los derechos civiles o políticos frustra el ejercicio de los derechos económicos, sociales, culturales y ambientales, y viceversa.

Sin embargo, si se analiza la estructura axiológica del sistema, la \textbf{Dignidad Humana} constituye el núcleo fundacional, principio matriz y derecho sustantivo supremo del cual dimanan todos los demás derechos fundamentales \citep{scjnDignidad2016}. En armonía con la dignidad, el \textbf{derecho a la vida} opera como presupuesto fáctico y ontológico indispensable, pues sin la existencia biológica y digna de la persona ningún otro derecho puede ser ejercido. En síntesis, más que hablar de un derecho ``más importante'' en sentido excluyente, la dignidad de la persona opera como la premisa axiológica inderogable que da sentido y coherencia a la totalidad del catálogo iusfundamental.

\medskip
\textbf{3. ¿Cuál es la relación de los derechos humanos con el Estado de Derecho?}

La relación entre los derechos humanos y el Estado de Derecho es de \textbf{co-originariedad, recíproca necesidad y justificación sustantiva}. 

Históricamente, el Estado de Derecho surgió en su vertiente formal o legalista como el ``gobierno de las leyes'' frente al arbitrio de los gobernantes (\emph{rule of law}). No obstante, en la teoría constitucional moderna (el Estado constitucional y democrático de Derecho), el Estado no se legitima meramente por el respeto a los procedimientos formales de creación normativa, sino por su contenido sustantivo: \textbf{un Estado sólo es verdaderamente de Derecho en la medida en que somete todo el ejercicio de sus poderes a la protección y garantía efectiva de los derechos humanos} \citep{carbonell2019}.

Los derechos humanos operan simultáneamente como:
\begin{enumerate}[leftmargin=*,topsep=2pt,itemsep=1pt]
\item \textbf{Límite negativo infranqueable} al poder estatal, vedando cualquier actuación arbitraria que lesione la esfera de autonomía de la persona (la ``esfera de lo indecidible'' de Ferrajoli).
\item \textbf{Mandato positivo y directriz funcional} para las políticas públicas y la administración de justicia, obligando al Estado a remover los obstáculos fácticos que impiden la igualdad real.
\item \textbf{Criterio de legitimidad política y democrática}, puesto que la soberanía popular no puede transgredir los derechos fundamentales sin desvirtuar su propia naturaleza democrática.
\end{enumerate}
\end{forobox}

\section{Conclusiones y alcance dialógico}
El debate académico en foros virtuales no constituye un mero trámite administrativo de acumulación de participaciones, sino un ejercicio dialógico deliberativo en el que se contrasta la teoría del Estado con la eficacia material de las garantías constitucionales. En la cultura jurídica moderna, comprender la co-originariedad entre el Estado de Derecho y los derechos humanos es indispensable para trascender el legalismo formalista y afianzar la primacía de la persona sobre las estructuras del poder público.

\textbf{Alcance de réplicas:} Conforme a la consigna de la Semana 6, la interacción deliberativa se mantendrá atenta a las intervenciones que formulen los compañeros de grupo en el aula virtual; las réplicas pertinentes serán elaboradas de manera contextualizada una vez que se cuente con intervenciones publicadas por los pares.
\hypertarget{cuerpo-fin}{}

\vspace{10pt}
\begingroup\fontsize{10}{11.5}\selectfont
\setlength{\bibsep}{4pt}
\phantomsection
\begin{thebibliography}{10}
\bibitem[Bobbio(1991)]{bobbio1991}
Bobbio, N. (1991). \emph{El tiempo de los derechos}. Editorial Sistema.

\bibitem[Carbonell(2019)]{carbonell2019}
Carbonell, M. (2019). \emph{Teor{\'i}a de los derechos fundamentales y justicia constitucional}. Editorial Trotta.

\bibitem[C{\'a}mara de Diputados(2026)]{cpeum2026S2}
C{\'a}mara de Diputados del H. Congreso de la Uni{\'o}n. (2026). \emph{Constituci{\'o}n Pol{\'i}tica de los Estados Unidos Mexicanos} (reforma DOF 7 de octubre de 2026). \url{https://www.diputados.gob.mx/LeyesBiblio/pdf/CPEUM.pdf}

\bibitem[Ferrajoli(2001)]{ferrajoli2001}
Ferrajoli, L. (2001). \emph{Derechos y garant{\'i}as: La ley del m{\'a}s d{\'e}bil}. Editorial Trotta.

\bibitem[SCJN(2014)]{scjnParametro2014}
Suprema Corte de Justicia de la Naci{\'o}n. (2014). Jurisprudencia P./J. 20/2014 (10a.): \emph{Derechos humanos contenidos en la Constituci{\'o}n y en los tratados internacionales. Constituyen el par{\'a}metro de control de regularidad constitucional}. Gaceta del Semanario Judicial de la Federaci{\'o}n, D{\'e}cima {\'E}poca, Libro 5, Tomo I, p. 202.

\bibitem[SCJN(2016)]{scjnDignidad2016}
Suprema Corte de Justicia de la Naci{\'o}n. (2016). Jurisprudencia 1a./J. 37/2016 (10a.): \emph{Dignidad humana. Constituye una norma jur{\'i}dica que consagra un derecho fundamental}. Gaceta del Semanario Judicial de la Federaci{\'o}n, D{\'e}cima {\'E}poca, Libro 33, Tomo II, p. 633.

\bibitem[UnADM(2026)]{unadmForoS6}
Universidad Abierta y a Distancia de M{\'e}xico. (2026). \emph{Consigna de actividades Semana 6: S6. Foro de Derechos Humanos} [Docente Titular Ver{\'o}nica Ortiz L{\'o}pez]. Aula virtual UnADM, Curso 3432.
\end{thebibliography}
\endgroup

\end{document}
